% P-15. Complete proof; the accompanying exact checker supplies the finite certificate.
\subsection{The support-nine seven-planes}
\begin{finding}{P-15}{a border-apolarity obstruction for three support-nine seven-planes}
Let $k$ be an algebraically closed field of characteristic zero.  Put
$B=C=k^{\{0,1,2\}\times\{0,1,2\}}$, and for row and column pairs
$I=\{r<s\}$, $J=\{c<d\}$ put
\[
 Q^{\epsilon}_{I,J}=
 x_{rc}\otimes x_{sd}+x_{sd}\otimes x_{rc}
 +\epsilon(x_{rd}\otimes x_{sc}+x_{sc}\otimes x_{rd}),
 \qquad \epsilon\in\{-1,1\}.
\]
Here the two displayed factors are $B,C$, rather than multiplication in a
symmetric algebra.  In the following two ordered lists the notation $01$
means the pair $\{0,1\}$:
\begin{align*}
 \mathcal S_A&=((01,01),(01,02),(01,12),(02,01),(02,02),(02,12),(12,01)),\\
 \mathcal S_B&=((01,01),(01,02),(01,12),(02,01),(02,02),(12,01),(12,12)).
\end{align*}
For $A=k^7$ with basis $a_0,\ldots,a_6$, define
$T^{\epsilon}_{\mathcal S}=\sum_{q=0}^6a_q\otimes Q^{\epsilon}_{\mathcal S(q)}$.
Then
\[
 \underline R(T^{-}_{\mathcal S_A})\ge14,\qquad
 \underline R(T^{+}_{\mathcal S_A})\ge14,\qquad
 \underline R(T^{+}_{\mathcal S_B})\ge14.
\]
The first of these is the determinant coordinate type
$[0,1,3,6,7,9,18]$ in the four-by-four rectangle indexing; the other two
are the permanent types $[0,1,3,6,7,9,18]$ and $[0,1,3,6,7,18,21]$.
\end{finding}

\paragraph{Proof and finite certificate.}
We give all reductions, including the meaning and completeness of the
finite linear-arrangement enumeration.  The attached checker
\texttt{check\_support9\_sevenplanes.py} reconstructs the matrices below
from the displayed lists and signs.  It reads no project data and imports
no other project code.  Its ranks and kernel computations are over
$\mathbb Q$; there are no floating-point or finite-field rank decisions.
The three resulting verification records are
\texttt{support9\_det\_verification.json},
\texttt{support9\_perA\_verification.json}, and
\texttt{support9\_perB\_verification.json}.

\emph{The necessary extension.}
Let $U=T(C^*)\subset A\otimes B$.  The nine contractions are independent,
so $\dim U=9$.  The proper-incidence argument used in R04, equivalently
the degree $(2,1,0)$ and $(1,2,0)$ border-apolarity necessary conditions,
says that $\underline R(T)\le13$ would provide a $13$-plane
$E\supset U$ such that
\begin{align}
 \operatorname{rank}\Phi_A(E)&\le78,&
 \Phi_A(E):A\otimes E&\longrightarrow\Lambda^2A\otimes B,\label{eq:P15-A}\\
 \operatorname{rank}\Phi_B(E)&\le104,&
 \Phi_B(E):B\otimes E&\longrightarrow A\otimes\Lambda^2B.\label{eq:P15-B}
\end{align}
On a pure vector $a\otimes b\in E$, the two maps are
$a'\otimes(a\otimes b)\mapsto a\wedge a'\otimes b$ and
$b'\otimes(a\otimes b)\mapsto a\otimes(b\wedge b')$.
The bounds are $(7-1)13$ and $(9-1)13$.
This invocation is a necessary condition, not an assertion that these two
conditions alone characterize border rank.

The exact baseline matrices have dimensions and ranks
\[
\begin{array}{c|c|c|c}
 &\text{matrix on }U&\text{rank on }U&\text{allowed rank increment}\\\hline
 \Phi_A&189\times63&62&16\\
 \Phi_B&252\times81&80&24
\end{array}.
\]
These numbers are the same in the three cases.  The checker constructs
each entry by the displayed wedge formulas and verifies both ranks by
exact Gaussian elimination.

\emph{Reduction to weight subspaces.}
Let $(\mathbb G_m)^6$ act with weights $e_r+f_c$ on both $x_{rc}$ factors,
and weight $-(e_r+e_s+f_c+f_d)$ on the corresponding $a_q$.
It fixes $T$ and preserves $U$.  The closed rank conditions
\eqref{eq:P15-A}--\eqref{eq:P15-B}, together with $U\subset E$, form a
closed torus-invariant subvariety of $\operatorname{Gr}(13,A\otimes B)$.
Choose an integral cocharacter separating all distinct torus characters
on $A\otimes B$.  A limit at zero of its orbit exists by properness of
the Grassmannian.  The limit still satisfies all conditions, and its
eigenspace decomposition for this cocharacter is the torus weight
decomposition.  Consequently an admissible $E$ may be assumed torus
invariant.

Write $V=(A\otimes B)/U=\bigoplus_wV_w$.  For each weight $w$ the
ambient coordinate block consists of the pairs $(q,v)$ with
\[
 w=\operatorname{wt}(x_v)-\operatorname{wt}(Q_{\mathcal S(q)}).
\]
The image $U$ has one dimension in weight $-\operatorname{wt}(x_v)$ for
each of the nine variables, and zero dimensions in the other weights.
In an image-weight block, omit its final coordinate in lexicographic
$(q,v)$ order; the coefficient of the omitted coordinate in the image
vector is nonzero.  The remaining coordinates give a specified basis of
$V_w$.  In a nonimage block take every coordinate.  Thus no choice of
quotient representatives is left unspecified.

For $\mathcal S_A$, there are $33$ nonzero quotient blocks: $15$ of
dimension one, $15$ of dimension two, and three of dimension three.
For $\mathcal S_B$, the corresponding numbers are $34=16+16+2$.
The total quotient dimension is $54$ in both cases.
An admissible extension is therefore specified by subspaces
$F_w\subset V_w$ with $\sum_w\dim F_w=4$.

\emph{The rank formula in quotient weight blocks.}
For either orientation $o\in\{A,B\}$, let $K_o$ be the image of the
baseline map on $U$.  Quotient each codomain weight block by its part of
$K_o$.  In the checker this quotient is represented by an integer-row
basis $H$ for the left nullspace of the baseline block $M_0$.  It checks
$HM_0=0$, $\operatorname{rank}H=\#\operatorname{rows}H$, and
$\#\operatorname{rows}H=\#\operatorname{rows}M_0-\operatorname{rank}M_0$.
Thus $\ker H=\operatorname{im}M_0$, and this is an exact quotient map.
For a factor basis vector with weight $\lambda_j$, write
\[
 L^o_{w,j}:V_w\longrightarrow
       (\operatorname{codomain}\Phi_o/K_o)_{w+\lambda_j}
\]
for the resulting linear map.  Its matrix is $H$ times the explicitly
constructed wedge columns.  For any chosen $F_w$ the exact rank increment
is
\begin{equation}
 \Delta_o(E)=
 \sum_{\gamma}\dim\left(
       \sum_{w+\lambda_j=\gamma}L^o_{w,j}(F_w)\right).
 \label{eq:P15-increment}
\end{equation}
Different target weights form a direct sum.  Formula
\eqref{eq:P15-increment} proves both the weight-count bounds used below
and the validity of subsequently quotienting by already fixed components.

\emph{Initial saturation.}
A support is a set $S$ of at most four quotient weights whose capacities
sum to at least four.  There are $46\,119$ eligible supports for
$\mathcal S_A$ and $51\,985$ for $\mathcal S_B$.
Replace all $F_w$, for $w\in S$, by the whole $V_w$, obtaining a maximal
space $E_S$.  If an admissible $E$ has support $S$, then the kernels of
its wedge maps inject into those for $E_S$.  Equations
\eqref{eq:P15-A}--\eqref{eq:P15-B} require both kernels for $E$ to have
dimension at least $13$, so a kernel of dimension below $13$ for $E_S$
excludes $S$.
The full wedge matrices on $A\otimes B$ have sizes $189\times441$ and
$252\times567$.  Decomposing them into weight blocks and precomputing
the ranks for every local support of at most four weights makes this
finite saturation calculation small.  Every support of size at most two
is excluded.  The surviving supports, before further tests, are
\[
\begin{array}{c|r|r|r}
 &\text{eligible supports}&\text{three weights}&\text{four weights}\\\hline
 T^{-}_{\mathcal S_A}&46\,119&236&13\,019\\
 T^{+}_{\mathcal S_A}&46\,119&244&13\,518\\
 T^{+}_{\mathcal S_B}&51\,985&200&11\,783
\end{array}.
\]
There are only two possible remaining multiplicity patterns:
$(1,1,1,1)$ and $(2,1,1)$.

\emph{Completeness of the linear-arrangement enumeration.}
For linear maps $C_1,\ldots,C_N$ on $k^d$, start with the zero equation
space and repeatedly add the row space of one $C_j$, retaining each
result of rank below $d$ and identifying equal row spaces by exact
row reduction.  This enumerates every possible intersection
$K=\bigcap_{j\in Z}\ker C_j\ne0$.
To such a $K$ assign the indices for which $C_j|_K=0$.
Its actual stratum consists of the nonzero vectors of $K$ on which every
other $C_j$ is nonzero.  It is nonempty: each prohibited additional zero
locus is a proper linear subspace of $K$, and a finite union of proper
linear subspaces cannot cover a nonzero vector space over the infinite
field $k$.
Conversely, the exact zero pattern of any nonzero vector occurs in this
list, by intersecting all the kernels that contain that vector.  Thus
these are locally closed strata; they are not being replaced silently
by their linear spans.

If a subsequent quotient produces new kernel conditions inside $K$, an
intersection $K'\subset K$ is compatible with the original stratum
exactly when none of the originally required nonzero maps vanishes
identically on $K'$.  The same finite-union argument proves both directions.
The checker explicitly makes this compatibility test whenever fixed
components have been quotiented out.

\emph{Four weight lines.}
For a line $F_w=kx$ the characters $w+\lambda_j$, as $j$ varies in one
orientation, are distinct.  Hence its contribution has one dimension
precisely for those $j$ with $L^o_{w,j}x\ne0$.
The combined $7+9$ column-vanishing conditions are linear equations in
at most three coordinates.  The arrangement procedure above therefore
lists every exact joint zero pattern in both orientations.
For several lines, the number of distinct active output weights is a
lower bound for \eqref{eq:P15-increment}.
A preliminary bound sums each line's minimum increment and subtracts
the largest possible loss from repeated output weights; this is valid
because each column has rank at most one.

If an arrangement stratum has a one-dimensional kernel, its line is
fixed.  Compute the actual sum of the images of all fixed lines, and
exclude the case if its ranks already exceed $16$ or $24$.
Otherwise quotient by those fixed images.  In every case that survives
this step there are one or two parameter lines.
For one line, apply the arrangement procedure to its columns in the
new quotient, with the compatibility test just proved.  The rank
increment equals the number of its nonzero columns.
For two lines, do this for each line and use the number of distinct
active weights in their union.  Every resulting case is excluded.
Here are the complete counts; all ranks in the fixed-image step are
computed over $\mathbb Q$:
\[
\begin{array}{c|r|r|r}
 &T^-_{\mathcal S_A}&T^+_{\mathcal S_A}&T^+_{\mathcal S_B}\\\hline
 \text{supports excluded by minimum/overlap}&12\,307&12\,798&11\,519\\
 \text{further supports excluded by joint weight unions}&176&180&36\\
 \text{joint four-line states remaining}&7\,225&7\,697&2\,182\\
 \text{states excluded by fixed-image ranks}&6\,619&7\,057&1\,972\\
 \text{one-parameter-line states, all excluded}&592&626&210\\
 \text{two-parameter-line states, all excluded}&14&14&0\\
 \text{refined two-line state pairs, all excluded}&32&22&0
\end{array}.
\]

\emph{A double component and two lines.}
For a two-dimensional quotient block, its double component is fixed.
For a three-dimensional block, a two-plane has the unique form
$\ker n$ for a nonzero dual normal $n$, modulo scaling.
Let $L:k^3\to D$ be any one of its column maps.
The complete rank-drop rule is
\[
\begin{array}{c|c|l}
 \operatorname{rank}L&\operatorname{rank}(L|_{\ker n})&\text{condition}\\\hline
 0&0&\text{always}\\
 1&0&n\text{ spans the row space of }L\\
 1&1&\text{otherwise}\\
 2&1&n\text{ annihilates }\ker L\\
 2&2&\text{otherwise}\\
 3&2&\text{always}.
\end{array}
\]
Indeed, in rank one the restriction is zero exactly when
$\ker n=\ker L$; in rank two the one-dimensional kernel of $L$ lies in
$\ker n$ exactly in the stated case; the other two rows are immediate.
Both nontrivial drop conditions are linear equations in $n$.
For rank one, equations cutting out the line $\operatorname{row}L$ are
obtained from a basis of its annihilator.  For rank two, the equation is
evaluation on a generator of $\ker L$.

The same arrangement algorithm therefore lists every possible rank
pattern of the double component, jointly in the two orientations.
Combine it with the two line patterns.  At each output weight, the
dimension of the sum in \eqref{eq:P15-increment} is at least the largest
of the dimensions contributed by the three components.  Summing these
maxima excludes every case:
\[
\begin{array}{c|r|r}
 &\text{arrangement states tested}&\text{states excluded}\\\hline
 T^-_{\mathcal S_A}&42\,274&42\,274\\
 T^+_{\mathcal S_A}&45\,489&45\,489\\
 T^+_{\mathcal S_B}&31\,796&31\,796.
\end{array}
\]
This completes the enumeration of every torus-fixed four-dimensional
extension quotient.  There is no admissible $E$, contradicting the
necessary condition for border rank at most $13$.  The three asserted
lower bounds follow.\qed

\paragraph{Algebraic-geometric dictionary for P-15.}
The tensor and both wedge maps are maps between finite free
$\mathbb Z$-modules.  Their rank conditions are determinantal closed
subschemes of the relevant relative Grassmannian.  The incidence
condition $U\subset E$ is a closed condition on the tautological bundle;
the Grassmannian's projectivity gives the proper projection used in R04
and the cocharacter specialization used here.  The acting split torus is
the group scheme $(\mathbb G_m)^6$, and its character decomposition is
an integral representation-ring decomposition.  Exterior powers in the
two wedge maps are the usual $\lambda$-operations on locally free sheaves.
The locally closed arrangement strata are obtained by intersections of
linear closed subschemes and complements of further such subschemes.
These are precisely the scheme, group-scheme, and $K_0$ constructions
identified in the paper's common Grothendieck locator table.
The certificate uses characteristic-zero rational linear algebra;
no statement about bad primes, purity, an $\ell$-adic weight filtration,
or a motivic realization is inferred from it.
